Language models can identify rhyming words despite never hearing speech, raising the question: where is this phonetic knowledge encoded?
We investigate the location and geometry of phonetic representations in \gptmed (24 layers, 1024 dimensions) using 13{,}105 single-token English words with IPA transcriptions from \wikipron.
We train per-phoneme linear probes at every layer of the residual stream and find that phonetic information is \emph{front-loaded}: probe accuracy peaks at layer~0 (macro F1 = 0.41) and does not improve in deeper layers, contrasting sharply with semantic features that build through the network.
We then analyze phoneme direction vectors---mean activation differences for words containing versus not containing each phoneme---and find that each phoneme has a distinct direction (mean pairwise $|\text{cosine similarity}| = 0.13$, $5.3\times$ above random), but these directions are not orthogonal subspaces; they become increasingly aligned in deeper layers as semantic representations develop.
Finally, we test context-dependent rhyming with \gptfour on 12 heteronyms and find a systematic primary-pronunciation bias: 92\% accuracy on primary but only 58\% on secondary pronunciations.
Together, these results show that rhyming knowledge lives primarily in the token embedding table, is carried but not enriched through the network, and that resolving context-dependent pronunciation remains an open challenge for current models.
